\subsection{EXISTENCE OF INTEGRALS OF A LINEAR DIFFERENTIAL EQUATION}

Let E be a complete normed space over the field R, and J an interval in R, not reducing to a point. One says that the differential equation

\[
{\frac {d \mathbf {x}}{d t}} = \mathbf {f} (t, \mathbf {x})\tag{1}
\]

where \(\mathbf{f}\) is defined on \(\mathbf{J} \times \mathbf{E}\) , is a linear equation if, for every \(t \in \mathbf{J}\) , the map \(\mathbf{x} \mapsto \mathbf{f}(t, \mathbf{x})\) is a continuous affine linear map \(^1\) from E into itself; if one puts \(\mathbf{b}(t) = \mathbf{f}(t, 0)\) the map \(\mathbf{x} \mapsto \mathbf{f}(t, \mathbf{x}) - \mathbf{f}(t, 0) = \mathbf{f}(t, \mathbf{x}) - \mathbf{b}(t)\) is then a continuous linear map from E to itself; from now on we shall denote this map by \(A(t)\) and write \(A(t).\mathbf{x}\) , (or simply \(A(t)\mathbf{x}\) ) for its value at the point \(\mathbf{x} \in E\) ; thus the linear differential equation (1) may be written

\[
{\frac {d \mathbf {x}}{d t}} = A (t). \mathbf {x} + \mathbf {b} (t)\tag{2}
\]

where \(\mathbf{b}\) is a map from J into E; when \(\mathbf{b} = 0\) one says that the linear differential equation (2) is homogeneous.

Examples. 1) When E is of finite dimension n over R one can identify the endomorphism \(A(t)\) with its matrix \(\left(a_{ij}(t)\right)\) with respect to any basis of E (Alg., II, p.~343); when one identifies a vector \(x \in E\) with the column matrix \(\left(x_{j}\right)\) of its components with respect to the basis of E under consideration the expression \(A(t)\) .x conforms to the general conventions of Algebra (Alg., II, p.~343, prop. 2). In this case, equation (2) is equivalent to the system of scalar differential equations

\[
\frac {d x _ {i}}{d t} = \sum_ {j = 1} ^ {n} a _ {i j} (t) x _ {j} + b _ {i} (t) \quad (1 \leqslant i \leqslant n).\tag{3}
\]

\begin{enumerate}
\def\labelenumi{\arabic{enumi})}
\setcounter{enumi}{1}
\tightlist
\item
  Let \(\mathbf{G}\) be a complete normed algebra over \(\mathbf{R}\) , and \(\mathbf{a}(t)\) , \(\mathbf{b}(t)\) and \(\mathbf{c}(t)\) three maps from \(\mathbf{J}\) into \(\mathbf{G}\) ; the equation
\end{enumerate}

\[
\frac {d \mathbf {x}}{d t} = \mathbf {a} (t) \mathbf {x} + \mathbf {x b} (t) + \mathbf {c} (t)
\]

is a linear differential equation; here \(A(t)\) is the linear map \(\mathbf{x} \mapsto \mathbf{a}(t)\mathbf{x} + \mathbf{x}\mathbf{b}(t)\) of \(G\) to itself.

For every \(t \in J\) , \(A(t)\) is an element of the set \(\mathcal{L}(\mathrm{E})\) of continuous linear maps from E to itself (continuous endomorphisms of E); one knows (Gen.~Top., X, p.~298) that \(\mathcal{L}(\mathrm{E})\) , endowed with the norm \(\|U\| = \sup_{\|x\| \leqslant 1} \|Ux\|\) is a complete normed algebra over the field R, and that \(\|UV\| \leqslant \|U\|\) \(\|V\|\) .

Throughout this section we shall assume that the following conditions are satisfied:

\begin{enumerate}
\def\labelenumi{\alph{enumi})}
\item
  The map \(t \mapsto A(t)\) of \(\mathbf{J}\) into \(\mathcal{L}(\mathrm{E})\) is regulated.
\item
  The map \(t \mapsto \mathbf{b}(t)\) of J into E is regulated.
\end{enumerate}

When \(\mathrm{E}\) has dimension \(n\) , \(\mathcal{L}(\mathrm{E})\) is isomorphic to \(\mathbf{R}^{n^2}\) (as a topological vector space) and condition \(a\) ) means that each of the elements \(a_{ij}(t)\) of the matrix \(A(t)\) is a regulated function on \(\mathrm{J}\) .

Since \(\left\| A(t')\mathbf{x} - A(t)\mathbf{x}\right\| \leqslant \left\| A(t') - A(t)\right\| \| \mathbf{x}\|\) , the map

\[
t \mapsto A (t). \mathbf {x} + \mathbf {b} (t)
\]

is regulated for every \(\mathbf{x} \in \mathrm{E}\) ; further,

\[
\left\| A (t) \mathbf {x} _ {1} - A (t) \mathbf {x} _ {2} \right\| = \left\| A (t) \left(\mathbf {x} _ {1} - \mathbf {x} _ {2}\right) \right\| \leqslant \left\| A (t) \right\| \left\| \mathbf {x} _ {1} - \mathbf {x} _ {2} \right\|
\]

for any \(t \in J\) and \(x_{1}\) , \(x_{2}\) in E; in other words, the right-hand side of (2) satisfies the conditions of lemma 1 of IV, p.~165 and is Lipschitz with respect to the regulated function \(\|A(t)\|\) on \(J \times E\) . In consequence (IV, p.~173, cor. 2):

THEOREM 1. Let \(t \mapsto A(t)\) be a regulated map of \(\mathbf{J}\) into \(\mathcal{L}(\mathrm{E})\) , and \(t \mapsto \mathbf{b}(t)\) be a regulated map of \(\mathbf{J}\) into \(\mathrm{E}\) . For every point \((t_0, \mathbf{x}_0)\) of \(\mathbf{J} \times \mathbf{E}\) the linear equation (2) admits one and only one solution defined on all of \(\mathbf{J}\) and equal to \(\mathbf{x}_0\) at the point \(t_0\) .

